% In this file you should put macros to be used only by
% the printed version. Of course they should have a corresponding
% version in macros/web.tex.
% Typically the printed version could have more fancy decorations.
% This should be a very short file.
%
% This file starts with dummy macros that ensure the pdf
% compiler will ignore macros provided by plasTeX that make
% sense only for the web version, such as dependency graph
% macros.


% Dummy macros that make sense only for web version.
\newcommand{\lean}[1]{}
\newcommand{\discussion}[1]{}
\newcommand{\leanok}{}
\newcommand{\mathlibok}{}
\newcommand{\notready}{}
% Make sure that arguments of \uses and \proves are real labels, by using invisible refs:
% latex prints a warning if the label is not defined, but nothing is shown in the pdf file.
% It uses LaTeX3 programming, this is why we use the expl3 package.
\ExplSyntaxOn
\NewDocumentCommand{\uses}{m}
 {\clist_map_inline:nn{#1}{\vphantom{\ref{##1}}}%
  \ignorespaces}
\NewDocumentCommand{\proves}{m}
 {\clist_map_inline:nn{#1}{\vphantom{\ref{##1}}}%
  \ignorespaces}
\ExplSyntaxOff

% ---------------------------------------------------------------------------
% cleveref names for the shared-counter theorem environments.
%
% macros/common.tex declares proposition/lemma/corollary/definition as
% `\newtheorem{...}[theorem]{...}` so that they all share one numbering
% sequence. cleveref names a reference after the *counter* the environment
% steps, not after the environment, so all four rendered as "theorem N" --
% `\cref{def:pauli-group-element}` on a definition printed "theorem 2".
%
% aliascnt fixes this without disturbing the numbering: `\newaliascnt{definition}
% {theorem}` makes \c@definition and \c@theorem the very same counter, so
% re-declaring the environment against `definition` keeps the shared sequence
% while giving cleveref a distinct counter name to look up. The environments
% already exist at this point, so \let each one back to \relax first --
% \newtheorem refuses to redefine a live command.
%
% Print-only: plasTeX takes the name from the environment's caption rather than
% from a counter, so the web version already rendered these correctly.
% ---------------------------------------------------------------------------
\usepackage{aliascnt}

\newcommand{\bpaliascounter}[2]{%
  \newaliascnt{#1}{theorem}%
  \expandafter\let\csname #1\endcsname\relax
  \expandafter\let\csname end#1\endcsname\relax
  \newtheorem{#1}[#1]{#2}%
  \aliascntresetthe{#1}%
}
\theoremstyle{plain}
\bpaliascounter{proposition}{Proposition}
\bpaliascounter{lemma}{Lemma}
\bpaliascounter{corollary}{Corollary}
\theoremstyle{definition}
\bpaliascounter{definition}{Definition}

% ---------------------------------------------------------------------------
% Bracket-safe theorem notes.
%
% amsthm grabs the optional note of \begin{theorem}[...] with a scan delimited
% by the first `]`, so a title containing brackets -- every `[[n,k,d]]` title in
% this blueprint -- ends the note early and drops the remainder into math mode.
% Six nodes hit this (def:steane7, def:shor9, def:five-qubit, def:four-qubit,
% def:iceberg, thm:base-distance).
%
% Re-grab the note with an xparse `o` argument, which balances nested brackets,
% and hand it to the original environment wrapped in braces so amsthm's scan
% never sees a bracket. Only the pdf build needs this: plasTeX parses the
% optional argument itself and renders these titles correctly as they stand.
% ---------------------------------------------------------------------------
\newcommand{\bracketsafethm}[1]{%
  \expandafter\let\csname bpsafe@#1\expandafter\endcsname\csname #1\endcsname
  \expandafter\let\csname endbpsafe@#1\expandafter\endcsname\csname end#1\endcsname
  \RenewDocumentEnvironment{#1}{o}
    {\IfValueTF{##1}{\csname bpsafe@#1\endcsname[{##1}]}{\csname bpsafe@#1\endcsname}}
    {\csname endbpsafe@#1\endcsname}%
}
\bracketsafethm{theorem}
\bracketsafethm{proposition}
\bracketsafethm{lemma}
\bracketsafethm{corollary}
\bracketsafethm{definition}
